% !TEX root = ../../main.tex

\section{Regularity tools}\label{sec:regularity-tools}

The limit varifold $V$ from the min-max construction is stationary (Proposition~\ref{thm:CLS-stationary}). To promote it to a smooth embedded minimal hypersurface, we need a regularity theorem. The classical tool is the Schoen--Simon regularity theorem~\cite{SS81}, which applies to stable codimension~$1$ integral varifolds $V$ that are either \emph{almost minimizing} in sufficiently small annuli or satisfy $\mathcal{H}^{n-2}(\sing V) = 0$; under either condition, $\sing V = \varnothing$ for $2 \leq n \leq 6$. In the Almgren--Pitts theory~\cite{Alm65, Pit81} and its refinements by Colding--de~Lellis~\cite{CL03} and de~Lellis--Tasnady~\cite{DLT13}, the almost minimizing property is produced by a combinatorial selection procedure, and Schoen--Simon then gives full regularity.

In the non-excessive framework of Chodosh--Liokumovich--Spolaor~\cite{CLS22}, however, the geometric condition is \emph{one-sided homotopic minimization} (Definition~\ref{def:homotopic-minimizer}), which provides stability but does not immediately yield the almost minimizing condition in annuli that Schoen--Simon requires. In~\cite{CLS22}, this gap is bridged by invoking the Almgren--Pitts regularity theory as a black box. To close this gap without relying on Almgren--Pitts, we turn to Wickramasekera's regularity theorem~\cite{Wic14}, which applies to a broader class of stable codimension~$1$ integral varifolds defined by three hypotheses. We state them following~\cite[Section~2]{Wic14}, adapted from the Euclidean ball $B_2^{n+1}(0)$ to a Riemannian manifold.

\begin{definition}[The class $\mathcal{S}_\alpha$~{\cite[Section~2]{Wic14}}]\label{defn:class-S-alpha}
Fix $\alpha \in (0, 1)$. Let $V$ be an integral $n$-varifold on an open subset $U$ of a smooth $(n+1)$-dimensional Riemannian manifold $(N, g)$, with $\|V\|(U) < \infty$. We say $V \in \mathcal{S}_\alpha$ if $V$ satisfies the following three conditions:
\begin{itemize}
\item[$(\mathcal{S}1)$] \textbf{Stationarity.} $V$ has vanishing first variation:
\[
    \delta V(X) = \int \operatorname{div}_S X \, dV(x, S) = 0 \quad \text{for every } X \in \mathfrak{X}_c(U).
\]

\item[$(\mathcal{S}2)$] \textbf{Stability.} For each open set $\Omega \subset U$ such that $\sing V \cap \Omega = \varnothing$ (if $2 \leq n \leq 6$) or $\mathcal{H}^{n-7+\gamma}(\sing V \cap \Omega) = 0$ for every $\gamma > 0$ (if $n \geq 7$),
\[
    \int_{\reg V \cap \Omega} |A|^2 \zeta^2 \, d\mathcal{H}^n \leq \int_{\reg V \cap \Omega} |\nabla \zeta|^2 \, d\mathcal{H}^n \quad \text{for all } \zeta \in C^1_c(\reg V \cap \Omega),
\]
where $A$ is the second fundamental form of $\reg V$. Equivalently, $V$ has non-negative second variation for normal deformations compactly supported in $\Omega \setminus \sing V$.

\item[$(\mathcal{S}3)$] \textbf{$\alpha$-structural hypothesis.}\label{defn:alpha-structural} For each $Z \in \sing V$, there exists \emph{no} $\rho > 0$ such that $\spt\|V\| \cap B_\rho(Z)$ is equal to the union of finitely many embedded $C^{1,\alpha}$ hypersurfaces-with-boundary of $B_\rho(Z)$, all having a common $C^{1,\alpha}$ boundary containing $Z$, with no two intersecting except along their common boundary.
\end{itemize}
\end{definition}

\begin{figure}[ht]
    \centering
    \begin{tikzpicture}[scale=0.85]
      % Edge Gamma: shared edge of three half-planes, with perspective
      \coordinate (A) at (-1.6,-0.5);
      \coordinate (B) at (1.6,0.5);

      % M_1 (upper half-plane): parallelogram above L
      \coordinate (A1) at ($(A)+(0.6,2.0)$);
      \coordinate (B1) at ($(B)+(0.6,2.0)$);

      % M_2 (lower-right half-plane): parallelogram to the right of L
      \coordinate (A2) at ($(A)+(2.0,-0.6)$);
      \coordinate (B2) at ($(B)+(2.0,-0.6)$);

      % M_3 (lower-left half-plane): parallelogram to the left of L
      \coordinate (A3) at ($(A)+(-1.8,-1.2)$);
      \coordinate (B3) at ($(B)+(-1.8,-1.2)$);

      % Intersection of A3--B3 with A--A2 (where M_3 passes behind M_2)
      \coordinate (I) at (-0.56,-0.81);

      % Fill and draw half-planes (back to front)
      % M_3 (behind): edges occluded by M_2 are dashed
      \fill[fillA, opacity=0.2] (A) -- (B) -- (B3) -- (A3) -- cycle;
      \draw[thick, bindA] (A) -- (A3);
      \draw[thick, bindA, dashed] (B) -- (B3);
      \draw[bindA] (A3) -- (I);
      \draw[bindA, dashed] (I) -- (B3);
      \node[bindA] at (-0.9,-1.15) {$M_3$};

      % M_1 (upper)
      \fill[fillA, opacity=0.2] (A) -- (B) -- (B1) -- (A1) -- cycle;
      \draw[thick, bindA] (A) -- (A1);
      \draw[thick, bindA] (B) -- (B1);
      \draw[bindA] (A1) -- (B1);
      \node[bindA] at (0.33,1.25) {$M_1$};

      % M_2 (lower-right)
      \fill[fillA, opacity=0.2] (A) -- (B) -- (B2) -- (A2) -- cycle;
      \draw[thick, bindA] (A) -- (A2);
      \draw[thick, bindA] (B) -- (B2);
      \draw[bindA] (A2) -- (B2);
      \node[bindA] at (1.1,-0.55) {$M_2$};

      % Edge Gamma (drawn on top)
      \draw[very thick] (A) -- (B);
      \node[above left] at ($(A)!0.28!(B)$) {$\Gamma$};

      % Singular point Z
      \coordinate (Z) at ($(A)!0.5!(B)$);
      \fill (Z) circle (1.5pt);
      \node[below] at ($(Z)+(0,-0.1)$) {\small $Z$};
    \end{tikzpicture}
    \caption{A configuration excluded by the $\alpha$-structural hypothesis $(\mathcal{S}3)$: three $C^{1,\alpha}$ hypersurfaces-with-boundary $M_1, M_2, M_3$ meeting along a common boundary $\Gamma$. The point $Z \in \Gamma$ is a singular point whose tangent cone is a union of three half-hyperplanes.}
    \label{fig:alpha-structural}
\end{figure}

\begin{remark}\label{rem:S-alpha-remarks}
\begin{enumerate}[label=\textup{(\roman*)}]
    \item Since $\reg V$ is open and dense in $\spt\|V\|$ by Allard's regularity theorem, $(\mathcal{S}2)$ is never vacuous; moreover $\mathcal{H}^{n-1}(\sing V) = 0$ trivially implies $(\mathcal{S}3)$, so $\mathcal{S}_\alpha$ contains all stable minimal hypersurfaces with small singular set~\cite[Remarks~(1)--(2) following~$(\mathcal{S}3)$]{Wic14}.
    \item For $V \in \mathcal{S}_\alpha$, hypothesis $(\mathcal{S}3)$ is equivalent to requiring that no tangent cone at a singular point is supported by \emph{three or more} distinct half-hyperplanes meeting along a common $(n-1)$-dimensional subspace; this equivalence relies on the Sheeting Theorem~\cite[Theorem~3.3]{Wic14} and the Minimum Distance Theorem~\cite[Theorem~3.4]{Wic14}; see also~\cite[Remark~(3) following~$(\mathcal{S}3)$]{Wic14}. The case of exactly two half-hyperplanes is already excluded by stationarity: a stationary cone supported on $N$ half-hyperplanes $P_1^+, \ldots, P_N^+$ meeting along a common edge $L$ satisfies the balancing condition $\sum_{j=1}^N \nu_j = 0$, where $\nu_j$ is the outward unit conormal of $P_j^+$ along $L$. When $N = 2$, this forces $\nu_1 = -\nu_2$, so the two half-hyperplanes form a single hyperplane and $Z$ is a regular point.
\end{enumerate}
\end{remark}

The main regularity result is:

\begin{theorem}[Regularity and Compactness~{\cite[Theorem~3.1, manifold version Theorem~6.1]{Wic14}}]\label{thm:wickramasekera}
    Let $(N^{n+1}, g)$ be a smooth Riemannian manifold and $\alpha \in (0, 1/2)$. If $V \in \mathcal{S}_\alpha$ on $N$ with $\|V\|(N) < \infty$, then:
    \begin{enumerate}[label=\textup{(\alph*)}]
        \item $\sing V = \varnothing$ if $2 \leq n \leq 6$;
        \item $\sing V$ is discrete if $n = 7$;
        \item $\mathcal{H}^{n-7+\gamma}(\sing V) = 0$ for each $\gamma > 0$ if $n \geq 8$.
    \end{enumerate}
    In particular, for $2 \leq n \leq 6$, $\spt\|V\|$ is a smooth embedded minimal hypersurface.
\end{theorem}

Wickramasekera's theorem strictly subsumes Schoen--Simon: if $\mathcal{H}^{n-2}(\sing V) = 0$, then $(\mathcal{S}3)$ holds vacuously (a junction of $C^{1,\alpha}$ hypersurfaces-with-boundary along an $(n-1)$-dimensional boundary would produce an $(n-1)$-dimensional singular set). Similarly, the almost minimizing property implies $\mathcal{H}^{n-2}(\sing V) = 0$ by Schoen--Simon, and thus implies $(\mathcal{S}3)$.

